\documentclass[11pt,a4paper]{article}
\usepackage[T1]{fontenc}
\usepackage[utf8]{inputenc}
\usepackage{amsmath,amssymb,amsthm}
\usepackage[margin=1in]{geometry}
\usepackage[colorlinks=true,linkcolor=blue,urlcolor=blue]{hyperref}
\theoremstyle{plain}
\newtheorem{theorem}{Theorem}
\newtheorem{lemma}[theorem]{Lemma}
\newtheorem{proposition}[theorem]{Proposition}
\newtheorem{corollary}[theorem]{Corollary}
\theoremstyle{definition}
\newtheorem{remark}[theorem]{Remark}

\title{Recovering the unwritten recurrences of the table OEIS A236819}
\author{Adrian Perez Fontelles\\ \small Independent researcher}
\date{2 September 2026}
\begin{document}
\maketitle

\begin{abstract}
OEIS A236819 is a two-dimensional table $T(n,k)$ whose empirical recurrences are, for several of
its lines, given only as an ORDER --- ``k=2: [order 34]'' and the like --- with the
coefficients in a linked file rather than in the entry. Those conjectures can still be settled,
and the recurrences recovered. A line of the table is a fixed-width or fixed-height array
count, hence a walk count on finitely many vertices, hence satisfies some monic recurrence of
order at most that number; Berlekamp--Massey on twice as many exact terms returns the MINIMAL
one. Where its order equals the order the entry states --- and where the same holds after the
first few terms are dropped --- any recurrence of that order the line satisfies has a
characteristic polynomial that is a multiple of the minimal one and of equal degree, hence is
that polynomial. This note settles column 2, column 3.
\end{abstract}

\noindent\small 2020 Mathematics Subject Classification. 05A15, 11B37, 15A18.\normalsize

\section{The table and the unwritten conjectures}

OEIS A236819 is ``T(n,k)=Number of (n+1)X(k+1) 0..2 arrays with the maximum plus the lower median minus the minimum of every 2X2 subblock equal.''. It is read by antidiagonals and begins
\[
81, 303, 303, 1162, 1791, 1162, 4627, 10702,\ \dots
\]
The entry carries, among its empirical claims:
\begin{quote}\small
k=2: [order 34]\\
k=3: [order 87]
\end{quote}
As of the ``Last modified'' line on the live entry (Jul 23 2025 09:35:45, revision 6) these are still
recorded as empirical, and the recurrences themselves are not in the entry text.

\section{Why a line of the table is a walk count}

Fix the index that the line fixes. The entry defines $T(n,k)$ as a number of arrays of a shape
determined by $n$ and $k$, subject to a condition local to a bounded window of consecutive
lines. Substituting the fixed index into the entry's own wording turns the definition into an
ordinary one-parameter array count, and for such a count the admissible windows are the
vertices of a finite digraph and an array is a walk. So the line is a walk count on $S$
vertices, and by the Cayley--Hamilton theorem it satisfies a monic integer recurrence of order
at most $S$.

\section{Recovering the recurrences}

\begin{lemma}\label{lem:bm}
A sequence known to satisfy some linear recurrence of order at most $S$ is determined, with its
minimal recurrence, by its first $2S$ terms; Berlekamp--Massey applied to them returns that
minimal recurrence exactly.
\end{lemma}

\subsection*{Column $k=2$}
Substituting the fixed index into the entry's wording gives an ordinary one-parameter array
count, a walk on $S=63$ vertices. Berlekamp--Massey on $2S=126$ exact terms
returns a minimal recurrence of order $34$ --- the order the entry states --- with
integer coefficients
\[
(c_1,\dots,c_{34})=(21,\ -100,\ -822,\ 9025,\ -6166,\ -227164,\ 761691,\ 2056841,\ -14904183,\ 4818303,\ 131300778,\ -239600971,\ -518603618,\ 1978489660,\ 49780225,\ -7775913759,\ 7725412959,\ 14487841224,\ -31018825865,\ -3698011866,\ 55220616016,\ -33860955124,\ -43168404248,\ 59680526360,\ 1187247872,\ -39820734688,\ 18284518864,\ 8552972960,\ -9006433472,\ 1006397632,\ 1171967872,\ -352288384,\ -23834624,\ 14026752).
\]
so that $a(m)=\sum_{i=1}^{34}c_i\,a(m-i)$ for every $m$. Removing the first
$34+4$ terms and repeating gives order $34$ again, which is what the
identification below needs.

\subsection*{Column $k=3$}
Substituting the fixed index into the entry's wording gives an ordinary one-parameter array
count, a walk on $S=164$ vertices. Berlekamp--Massey on $2S=328$ exact terms
returns a minimal recurrence of order $87$ --- the order the entry states --- with
integer coefficients
\[
(c_1,\dots,c_{87})=(30,\ -50,\ -6914,\ 56825,\ 620000,\ -8823890,\ -21231056,\ 712649613,\ -686140706,\ -36295497697,\ 112914575365,\ 1245544005845,\ -6347365002307,\ -28888854584644,\ 225736251950572,\ 407765750489979,\ -5738443521471890,\ -1149890832378082,\ 109370207550528170,\ -104043149918690319,\ -1594501820813188222,\ 3101906979159742404,\ 17804737602207351381,\ -53494269584001636377,\ -148703596391108404276,\ 665814537078362186550,\ 849882989649911503256,\ -6376552886165955266520,\ -2023985156801248518927,\ 48206244895647768402378,\ -19247842124293965032943,\ -290153664538467411399200,\ 296162804073970051128391,\ 1384408347307655156472189,\ -2299558632285427578387943,\ -5119547531344472937270184,\ 12733260417312019822688593,\ 13726171282244713700698824,\ -54417116335348544879521595,\ -20362079076182437439941352,\ 184498821186955151470086115,\ -25197955192795178476468826,\ -499980338491491519966112578,\ 283252270672941367733545224,\ 1075946560801602927700564035,\ -1070219996534568402363575063,\ -1793181179402408172139583460,\ 2738321572065985791249514692,\ 2159655119906908266891347986,\ -5286501216701410928180483672,\ -1435019391261447843334971234,\ 7945484919659038793673460710,\ -750953383754591097918759582,\ -9336452242143762257216928131,\ 3732624406398569798036613577,\ 8450925237534172928356779326,\ -5986149817181084701274288674,\ -5642029091504791094667419214,\ 6307064348820514803896774458,\ 2446334246314087416900936389,\ -4828081534818868150920326786,\ -295352484274418802990799580,\ 2739268892880168415854673272,\ -479676332431650123540919378,\ -1136711462917031055123804740,\ 436355264688675694178205640,\ 329135873666336160774199000,\ -206819120867888256993973816,\ -58352735087144698807931984,\ 63747597328776519191534160,\ 2920468939974786392062784,\ -13221788983829320255297824,\ 1370809733124892812403648,\ 1818519847589773335681280,\ -396461400258072719214080,\ -158322144308385044455936,\ 52657779601112847646720,\ 7823041651295994738688,\ -4024317062923539507200,\ -140605291014906247168,\ 179192664536189775872,\ -4762419016424751104,\ -4273680388290527232,\ 268993989203230720,\ 40627457116340224,\ -3611739520237568,\ 42044793815040).
\]
so that $a(m)=\sum_{i=1}^{87}c_i\,a(m-i)$ for every $m$. Removing the first
$87+4$ terms and repeating gives order $87$ again, which is what the
identification below needs.

\begin{theorem}
For each line above, the displayed recurrence holds for every index, and it is the recurrence
the entry records.
\end{theorem}

\begin{proof}
It holds by Lemma~\ref{lem:bm}. For the identification, the entry asserts that the line
satisfies SOME recurrence of the stated order, possibly only from some index on. Let $q$ be its
characteristic polynomial and $q_{\min}$ the minimal polynomial of the tail. Then
$q_{\min}\mid q$, and the second Berlekamp--Massey run --- on the line with its first few
terms removed --- shows $\deg q_{\min}$ equals the stated order, which is $\deg q$. Both are
monic, so $q=q_{\min}$.
\end{proof}

\section{Verification}

Three checks.

First, each line's model was compared against the entry's own published values for that line,
taken from the antidiagonal DATA read in either orientation and from the ``Table starts''
block, and agrees at every available term. This is the only point at which reading the entry's
English enters, and a misreading gives different counts.

Second, each recovered recurrence was evaluated directly on the model's values, with no
Berlekamp--Massey involved, and holds at every index tested.

Third, each order was computed twice by independent routes --- modulo a $61$-bit prime, where
every intermediate is a machine word, and again in exact rational arithmetic --- and once more
on the line with its first few terms dropped, which is what the identification requires. All
agree.

\begin{thebibliography}{9}
\bibitem{oeis} The OEIS Foundation, \emph{The On-Line Encyclopedia of Integer
Sequences}, \url{https://oeis.org}, 2026. Sequence A236819.
\bibitem{massey} J.~L.~Massey, Shift-register synthesis and BCH decoding, \emph{IEEE
Trans. Inform. Theory} \textbf{15} (1969), 122--127.
\bibitem{stanley} R.~P.~Stanley, \emph{Enumerative Combinatorics}, Volume 1, 2nd ed.,
Cambridge University Press, 2011. (Transfer-matrix method, Section 4.7.)
\bibitem{horn} R.~A.~Horn and C.~R.~Johnson, \emph{Matrix Analysis}, 2nd ed., Cambridge
University Press, 2013.
\end{thebibliography}

\end{document}
